\section*{Capítulo \ref{ch:g9:kinetic-energy-safety} --- La energía cinética y la seguridad vial}

\begin{solution}{exo:g9:kinetic-energy-safety:1}
El \omterm{def:g9:kinetic-energy-safety:joule}{julio} (\unit{J}): unos pocos \omterm{def:g9:kinetic-energy-safety:joule}{julios} suben este libro a
su estante; el corazón gasta alrededor de uno por latido. $E_k =
\frac12 m v^2$ --- $m$ en \omterm{def:g3:measuring-mass:units}{kilogramos}, $v$ en \omterm{def:g2:measuring-length:units}{metros} por
segundo y $E_k$ en \omterm{def:g9:kinetic-energy-safety:joule}{julios}.
\end{solution}

\begin{solution}{exo:g9:kinetic-energy-safety:2}
Liebre: $0.5 \times 2 \times 100 = \qty{100}{J}$. Jugador: $0.5 \times
80 \times 25 = \qty{1000}{J}$.
\end{solution}

\begin{solution}{exo:g9:kinetic-energy-safety:3}
Media \omterm{def:g3:measuring-mass:mass}{masa} más: factor $1.5$. Media \omterm{def:g7:motion-average-speed:formula}{velocidad} más: factor
$1.5^2 = 2.25$. Gana el mando de la \omterm{def:g7:motion-average-speed:formula}{velocidad} --- gana siempre,
porque lleva puesto el cuadrado.
\end{solution}

\begin{solution}{exo:g9:kinetic-energy-safety:4}
La \omterm{prop:g9:kinetic-energy-safety:stopping}{distancia de reacción} --- el viaje de un segundo, que crece en
proporción a $v$; la \omterm{prop:g9:kinetic-energy-safety:stopping}{distancia de frenado} --- la jubilación de
$E_k$, que crece con $v^2$.
\end{solution}

\begin{solution}{exo:g9:kinetic-energy-safety:5}
A \qty{50}{km/h}: $13$ de \qty{27}{m} --- alrededor de la mitad. A
\qty{130}{km/h}: $85$ de \qty{121}{m} --- un setenta por ciento. El
componente al cuadrado le crece por encima al proporcional: a
\omterm{def:g7:motion-average-speed:formula}{velocidad} alta, el frenado devora el total.
\end{solution}

\begin{solution}{exo:g9:kinetic-energy-safety:6}
A \omterm{def:g5:heat-and-insulation:heat}{calor} en los discos de freno --- el resplandor ámbar tras las
bajadas de montaña, el olor a frenos calientes. La
\omterm{def:g5:everyday-energy:energy}{energía} se pasa, no se borra nunca: la proposición de las cadenas
del capítulo infantil de la \omterm{def:g5:everyday-energy:energy}{energía}, todavía al mando.
\end{solution}

\begin{solution}{exo:g9:kinetic-energy-safety:7}
$36 \div 3.6 = \qty{10}{m}$; después \qty{20}{m}; después
\qty{30}{m} --- desfilando en proporción a la \omterm{def:g7:motion-average-speed:formula}{velocidad}, como
tiene que hacer un tramo de un segundo.
\end{solution}

\begin{solution}{exo:g9:kinetic-energy-safety:8}
Tu $E_k$ la fijan tu \omterm{def:g3:measuring-mass:mass}{masa} y la \omterm{def:g7:motion-average-speed:formula}{velocidad} del coche ---
ninguna correa la cambia. El cinturón civiliza la
\emph{jubilación}: para el cuerpo a lo largo de una distancia y un
tiempo mayores, con suavidad, en lugar de contra el salpicadero y de
golpe.
\end{solution}

\begin{solution}{exo:g9:kinetic-energy-safety:9}
El frenado escala con el cuadrado de la \omterm{def:g7:motion-average-speed:formula}{velocidad}: a
\qty{60}{km/h} (el doble), $8 \times 4 = \qty{32}{m}$; a
\qty{90}{km/h} (el triple), $8 \times 9 = \qty{72}{m}$.
\end{solution}

\begin{solution}{exo:g9:kinetic-energy-safety:10}
A \qty{30}{km/h}: reacción $\approx \qty{8.3}{m}$, frenado $13 \times
(30/50)^2 \approx \qty{4.7}{m}$ --- parada en unos \qty{13}{m}, frente
a \qty{27}{m} a cincuenta: menos de la mitad de carretera. \omterm{def:g5:everyday-energy:energy}{Energías}:
$(30/50)^2 = 0.36$ --- apenas un tercio de la \omterm{def:g5:everyday-energy:energy}{energía} de choque.
Las dos convencen; la línea que sirve para un cartel suele ser la de la
carretera --- ``a 30 te paras antes de llegar al niño; a 50 llegas
todavía en movimiento''.
\end{solution}

\begin{solution}{exo:g9:kinetic-energy-safety:11}
Dobla solo la parte del frenado --- la lluvia alarga la jubilación, no
el cerebro: a \qty{50}{km/h}, $14 + 26 = \qty{40}{m}$; a
\qty{90}{km/h}, $25 + 80 = \qty{105}{m}$. El segundo de reacción es
igual en seco que en mojado.
\end{solution}

\begin{solution}{exo:g9:kinetic-energy-safety:12}
Camión: $v = 25$ \unit{m/s}; $E_k = 0.5 \times 20000 \times 625 =
\qty{6.25e6}{J}$. Coche que lo iguala: $0.5 \times 1000 \times v^2 =
\num{6.25e6}$ da $v^2 = 12500$, $v \approx \qty{112}{m/s} \approx
\qty{400}{km/h}$ --- ningún coche de carretera se le acerca. Un camión
a \omterm{def:g7:motion-average-speed:formula}{velocidad} de autopista lleva \omterm{def:g5:everyday-energy:energy}{energías} de coche de carreras con
frenos de camión: cuando se le recalientan y le fallan en una bajada
larga, solo una rampa de grava puede jubilar megajulios con seguridad.
Los coches no la necesitan nunca; sus \omterm{def:g5:everyday-energy:energy}{energías} mueren en unos frenos
corrientes.
\end{solution}

\begin{solution}{pb:g9:kinetic-energy-safety:1}
\textbf{1.} $50 \div 3.6 = \qty{13.9}{m/s}$; $60 \div 3.6 =
\qty{16.7}{m/s}$.
\textbf{2.} \qty{13.9}{m} y \qty{16.7}{m}.
\textbf{3.} \qty{13}{m}; a sesenta, $13 \times (60/50)^2 = 13 \times
1.44 \approx \qty{18.7}{m}$.
\textbf{4.} Parada: $26.9$ frente a \qty{35.4}{m} --- el hueco del
cartel: unos \emph{ocho \omterm{def:g2:measuring-length:units}{metros} y medio} más, dos largos de coche
más allá de donde se ha parado quien iba a cincuenta.
\textbf{5.} A cincuenta: unos \qty{97}{kJ}. A noventa ($v =
\qty{25}{m/s}$): $0.5 \times 1000 \times 625 = \qty{312}{kJ}$.
\textbf{6.} $97 \div 9.8 \approx \qty{10}{m}$; $312 \div 9.8 \approx
\qty{32}{m}$.
\textbf{7.} Diez \omterm{def:g2:measuring-length:units}{metros}: una casa de tres plantas; treinta y dos: un
bloque de diez. Borrador: ``Un choque a 50 es una caída desde el tercer
piso. A 90, desde el décimo. Elige tu ventana''.
\textbf{8.} Las alturas miden con honradez la \omterm{def:g5:everyday-energy:energy}{energía} que
hay que jubilar --- ese número no lo cambia ninguna ingeniería. Las
zonas de deformación y los cinturones cambian la \emph{manera} de
jubilarla: \omterm{def:g2:measuring-length:units}{metros} de parada suave en vez de \omterm{def:g2:measuring-length:units}{centímetros} de
parada cruel --- por eso la analogía de la caída exagera las lesiones en
un coche moderno, y por eso la \omterm{def:g5:everyday-energy:energy}{energía} que retrata es real de
todas formas.
\textbf{9.} \qty{25}{m} a ciegas por segundo; una miradita de dos
segundos, \qty{50}{m} --- medio campo antes de que empiece siquiera el
segundo de reacción.
\textbf{10.} $65 + 25 = \qty{90}{m}$ --- el \omterm{def:g4:conductors-insulators:conductor}{conductor} cansado añade un
autobús y medio al total del \omterm{def:g4:conductors-insulators:conductor}{conductor} atento.
\textbf{11.} Se equivoca más a \omterm{def:g7:motion-average-speed:formula}{velocidades} \emph{bajas}: ahí la
parte del frenado, que va al cuadrado, es pequeña y el tramo de
reacción es casi toda la distancia de parada --- en ciudad, el segundo
\emph{es} el peligro. (A \omterm{def:g7:motion-average-speed:formula}{velocidad} alta domina el frenado --- pero
\qty{25}{m} por segundo tampoco son una broma allí.)
\textbf{12.} Por ejemplo: ``Diez más de \omterm{def:g7:motion-average-speed:formula}{velocidad}, diez
\omterm{def:g2:measuring-length:units}{metros} más de parada --- la \omterm{def:g7:motion-average-speed:formula}{velocidad} entra dos veces, una de
ellas al cuadrado. Un choque es una caída: la \omterm{def:g5:everyday-energy:energy}{energía} crece con
el cuadrado de la \omterm{def:g7:motion-average-speed:formula}{velocidad}. Un segundo de miradita son
veinticinco \omterm{def:g2:measuring-length:units}{metros} de ceguera --- la atención es la
\omterm{prop:g9:kinetic-energy-safety:stopping}{distancia de frenado} más corta de todas''.
\end{solution}
